\documentclass[12pt]{article}
\usepackage{amsmath,amssymb,amsfonts}
\usepackage[margin=1in]{geometry}

\begin{document}

\begin{center}
{\large\bfseries 5th Irish Mathematical Olympiad\\
2 May 1992}
\end{center}

\bigskip
\noindent\textbf{Paper 1}
\bigskip

\begin{enumerate}
\item Describe in geometric terms the set of points $(x,y)$ in the plane such that $x$ and $y$ satisfy the condition $t^2 + yt + x \ge 0$ for all $t$ with $-1 \le t \le 1$.

\item How many ordered triples $(x,y,z)$ of real numbers satisfy the system of equations \begin{align*} x^2 + y^2 + z^2 &= 9, \\ x^4 + y^4 + z^4 &= 33, \\ xyz &= -4\;? \end{align*}

\item Let $A$ be a nonempty set with $n$ elements. Find the number of ways of choosing a pair of subsets $(B, C)$ of $A$ such that $B$ is a nonempty subset of $C$.

\item In a triangle $ABC$, the points $A', B'$ and $C'$ on the sides opposite $A$, $B$ and $C$, respectively, are such that the lines $AA'$, $BB'$ and $CC'$ are concurrent. Prove that the diameter of the circumscribed circle of the triangle $ABC$ equals the product $|AB'|\cdot|BC'|\cdot|CA'|$ divided by the area of the triangle $A'B'C'$.

\item Let $ABC$ be a triangle such that the coordinates of the points $A$ and $B$ are rational numbers. Prove that the coordinates of $C$ are rational if, and only if, $\tan A$, $\tan B$ and $\tan C$, when defined, are all rational numbers.

\end{enumerate}

\newpage

\begin{center}
{\large\bfseries 5th Irish Mathematical Olympiad\\
2 May 1992}
\end{center}

\bigskip
\noindent\textbf{Paper 2}
\bigskip

\begin{enumerate}
\setcounter{enumi}{5}
\item Let $n > 2$ be an integer and let $m = \sum k^3$, where the sum is taken over all integers $k$ with $1 \le k < n$ that are relatively prime to $n$. Prove that $n$ divides $m$. (Note that two integers are \emph{relatively prime} if, and only if, their greatest common divisor equals 1.)

\item If $a_1$ is a positive integer, form the sequence $a_1, a_2, a_3, \ldots$ by letting $a_2$ be the product of the digits of $a_1$, etc.. If $a_k$ consists of a single digit, for some $k \ge 1$, $a_k$ is called a \emph{digital root} of $a_1$. It is easy to check that every positive integer has a unique digital root. (For example, if $a_1 = 24378$, then $a_2 = 1344$, $a_3 = 48$, $a_4 = 32$, $a_5 = 6$, and thus 6 is the digital root of 24378.) Prove that the digital root of a positive integer $n$ equals 1 if, and only if, all the digits of $n$ equal 1.

\item Let $a, b, c$ and $d$ be real numbers with $a \ne 0$. Prove that if all the roots of the cubic equation \[az^3 + bz^2 + cz + d = 0\] lie to the left of the imaginary axis in the complex plane, then \[ab > 0, \quad bc - ad > 0, \quad ad > 0.\]

\item A convex pentagon has the property that each of its diagonals cuts off a triangle of unit area. Find the area of the pentagon.

\item If, for $k = 1, 2, \ldots, n$, $a_k$ and $b_k$ are positive real numbers, prove that \[\sqrt[n]{a_1a_2\cdots a_n} + \sqrt[n]{b_1b_2\cdots b_n} \le \sqrt[n]{(a_1+b_1)(a_2+b_2)\cdots(a_n+b_n)}\] and that equality holds if, and only if, \[\frac{a_1}{b_1} = \frac{a_2}{b_2} = \cdots = \frac{a_n}{b_n}.\]

\end{enumerate}

\end{document}
